\documentclass{article}
\usepackage{amsmath, amssymb, amsthm}
\usepackage{hyperref}
\usepackage{geometry}
\geometry{margin=1in}

\title{Erd\H{o}s Problem \#826}
\date{Last updated: 2025-08-31}
\author{Source: erdosproblems.com}

\begin{document}
\maketitle

\noindent\textbf{Status:} OPEN \quad \textbf{No prize}

\noindent\textbf{Tags:} number theory

\medskip
\noindent\textbf{Problem Statement:}

Are there infinitely many $n$ such that, for all $k\geq 1$,\[\tau(n+k)\ll k?\]

\bigskip
\noindent\textbf{Remarks:}

A stronger form of [248].

Lau \cite{La26} has proved that there exists an absolute constant $C$ such that, for infinitely many $n$,\[\tau(n+k) \ll k^C\]for all $k\geq 1$.

\bigskip
\noindent\textbf{References:}

[La26] C. F. Lau, On the number of prime factors of consecutive integers. arXiv:2604.15042 (2026).

\bigskip
\noindent\small{Source: \url{https://www.erdosproblems.com/826}}
\end{document}
